\documentclass[11pt]{article}

\usepackage[margin=0.7in]{geometry}
\usepackage[T1]{fontenc}
\usepackage{lmodern}
\usepackage{microtype}
\usepackage{amsmath,amssymb}
\usepackage[dvipsnames]{xcolor}
\usepackage{enumitem}
\usepackage{needspace}
\usepackage[most]{tcolorbox}
\usepackage[hidelinks]{hyperref}

\hypersetup{
  pdftitle={Homework 4, MATH 327, Fall 2026},
  pdfauthor={Manish Patnaik},
  pdfsubject={MATH 327 homework problems}
}

\definecolor{courseblue}{HTML}{1F4E79}
\definecolor{lightblue}{HTML}{EAF2F8}
\definecolor{boxborder}{HTML}{B8CBDC}
\definecolor{darkgray}{HTML}{333333}
\definecolor{lightyellow}{HTML}{FFF9DC}
\definecolor{yellowborder}{HTML}{E4D58C}

\setlength{\parindent}{0pt}
\setlength{\parskip}{0.55em}
\setlist[enumerate]{leftmargin=2em, itemsep=0.25em, topsep=0.3em}

\newcommand{\Z}{\mathbb{Z}}
\newcommand{\R}{\mathbb{R}}
\newcommand{\C}{\mathbb{C}}
\newcommand{\F}{\mathbb{F}}
\newcommand{\SL}{\mathrm{SL}}
\newcommand{\GL}{\mathrm{GL}}
\newcommand{\SO}{\mathrm{SO}}
\newcommand{\Cl}{\operatorname{Cl}}
\newcommand{\Stab}{\operatorname{Stab}}

\newtcolorbox{exercise}[2][]{
  enhanced,
  colback=white,
  colframe=boxborder,
  colbacktitle=lightblue,
  coltitle=courseblue,
  fonttitle=\bfseries,
  title={Exercise #2},
  boxrule=0.6pt,
  arc=1mm,
  left=2.5mm,
  right=2.5mm,
  top=1.5mm,
  bottom=1.5mm,
  before skip=0.6em,
  after skip=0.6em,
  #1
}

\newtcolorbox{solution}{
  enhanced,
  breakable,
  colback=lightyellow,
  colframe=yellowborder,
  colbacktitle=lightyellow,
  coltitle=darkgray,
  fonttitle=\bfseries,
  title={Solution},
  boxrule=0.6pt,
  arc=1mm,
  left=2.5mm,
  right=2.5mm,
  top=1.5mm,
  bottom=1.5mm,
  before skip=0.2em,
  after skip=1em
}

\newcommand{\chapterheading}[1]{%
  \par\Needspace{12\baselineskip}\vspace{0.5em}
  {\color{courseblue}\Large\bfseries #1}\par\vspace{0.35em}
  {\color{boxborder}\hrule height 0.6pt}\vspace{0.35em}
}

\newcommand{\sectionheading}[1]{%
  \Needspace{8\baselineskip}
  \vspace{0.25em}
  {\color{darkgray}\large\bfseries #1}\par
}

\begin{document}

\begin{center}
  \colorbox{lightblue}{%
    \begin{minipage}{0.92\textwidth}
      \vspace{0.6em}
      \begin{tabular*}{\textwidth}{@{\extracolsep{\fill}}lr}
        {\color{courseblue}\bfseries\LARGE Homework 4} &
        {\color{courseblue}\bfseries\LARGE MATH 327}\\[0.4em]
        \multicolumn{2}{l}{\color{darkgray}\bfseries Fall 2026}\\[0.35em]
        \multicolumn{2}{l}{\color{darkgray}Covers Lectures 20--24}\\[0.2em]
        \multicolumn{2}{l}{\color{darkgray}Suggested completion date: Thursday, November 5}
      \end{tabular*}
      \vspace{0.5em}
    \end{minipage}}
\end{center}

\vspace{0.5em}

Unless otherwise indicated, the following exercises are from Michael Artin's
\emph{Algebra}, second edition. Exercises marked \emph{Supplementary} are
not from the text. This set is also good practice for Midterm~2.

% Problems drawn in part from Homework (2025)/2025-Winter-MATH327-Homework3.tex
% (Ch. 6: 9.2, 12.3; Ch. 7: 2.8(b), 3.2, 4.2, 5.6, and the SL_2(F_3) class equation).

%======================================================================
\chapterheading{Chapter 2: Groups}

\sectionheading{Section 12: Quotient Groups and the First Isomorphism Theorem}

\begin{exercise}{S.1 (Supplementary)}
Use the first isomorphism theorem to establish each of the following
isomorphisms. In each case, write down the homomorphism you use and
compute its kernel and image.
\begin{enumerate}[label=\textbf{(\alph*)}]
  \item $\C^{\times}/U\cong\R_{>0}$, where $U=\{z\in\C:|z|=1\}$ and
        $\R_{>0}$ is the group of positive reals under multiplication.
  \item $(\R,+)/\Z\cong U$.
  \item $S_4/V\cong S_3$, where $V=\{1,(12)(34),(13)(24),(14)(23)\}$.
        \textit{Hint.} Let $S_4$ act on the three ways of partitioning
        $\{1,2,3,4\}$ into two pairs,
        \[
          \pi_1=\{12\,|\,34\},\qquad
          \pi_2=\{13\,|\,24\},\qquad
          \pi_3=\{14\,|\,23\}.
        \]
\end{enumerate}
\end{exercise}

\begin{exercise}{S.2 (Supplementary)}
Let $N$ be a normal subgroup of a group $G$. The \emph{commutator} of
$a,b\in G$ is $[a,b]=aba^{-1}b^{-1}$.
\begin{enumerate}[label=\textbf{(\alph*)}]
  \item Prove that $G/N$ is abelian if and only if $[a,b]\in N$ for all
        $a,b\in G$.
  \item Deduce that for $n\geq 2$ every commutator in $S_n$ is an even permutation.
\end{enumerate}
\end{exercise}

%======================================================================
\chapterheading{Chapter 7: More Group Theory}

\sectionheading{Sections 2--3: The Class Equation and the Center}

\begin{exercise}{2.8(b)}
Determine the possible class equations of a nonabelian group of order $21$.
\end{exercise}

\begin{exercise}{3.2}
Let $Z$ be the center of a group $G$. Prove that if $G/Z$ is a cyclic
group, then $G$ is abelian, and hence $G=Z$.
\end{exercise}

\begin{exercise}{S.3 (Supplementary)}
Let $p$ be a prime.
\begin{enumerate}[label=\textbf{(\alph*)}]
  \item Let $G$ be a group of order $p^n$ with $n\geq1$. Use the class
        equation to prove that the center $Z(G)$ is nontrivial.
  \item Prove that every group of order $p^2$ is abelian.
\end{enumerate}
\end{exercise}

\begin{exercise}[breakable]{S.4 (Supplementary)}
Let $\F_3=\Z/3\Z$ and $G=\SL_2(\F_3)$, which has order $24$ (Lecture 22).
Put
\[
  U=\begin{pmatrix}1&1\\0&1\end{pmatrix},\qquad
  J=\begin{pmatrix}0&-1\\1&0\end{pmatrix}.
\]
\begin{enumerate}[label=\textbf{(\alph*)}]
  \item The group $G$ acts on the set of $8$ nonzero vectors in $\F_3^2$ by
        matrix multiplication. Show that this action is transitive, compute
        the stabilizer of $e_1=\binom10$, and use it to recover $|G|=24$.
  \item Show that $Z(G)=\{\pm I\}$.
  \item Compute the centralizers $C_G(U)$ and $C_G(J)$, and the sizes of
        the conjugacy classes of $U$ and $J$.
  \item Show that $U$ and $U^{-1}$ are \emph{not} conjugate in $G$.
        (They are conjugate in $\GL_2(\F_3)$.)
  \item Conclude that the class equation of $G$ is
        \[
          24=1+1+4+4+4+4+6,
        \]
        with class representatives $I$, $-I$, $U$, $U^{-1}$, $-U$, $-U^{-1}$, and $J$.
\end{enumerate}
\end{exercise}

\sectionheading{Sections 4--5: Conjugation in $S_n$ and $A_n$}

\begin{exercise}{5.6 (Adapted)}
\begin{enumerate}[label=\textbf{(\alph*)}]
  \item Find all subgroups of order $4$ of $S_4$, and decide which ones
        are normal.
  \item Use the class equation of $S_4$ to prove that $A_4$ is the only
        subgroup of $S_4$ of order $12$.
\end{enumerate}
\end{exercise}

\begin{exercise}{S.5 (Supplementary)}
\begin{enumerate}[label=\textbf{(\alph*)}]
  \item Use the splitting lemma from Lecture 23 to compute the class equation
        of $A_4$. Give a representative of each class, and decide whether
        $(123)$ and $(132)$ are conjugate in $A_4$.
  \item Deduce that $A_4$ has no subgroup of order $6$. (So the converse of
        Lagrange's theorem fails.)
\end{enumerate}
\end{exercise}

\begin{exercise}{4.2 (Adapted)}
Determine all normal subgroups of the symmetric group $S_5$.
\textit{Hint.} Use the class equation of $S_5$ from Lecture 23, or the
simplicity of $A_5$.
\end{exercise}

%======================================================================
\chapterheading{Chapter 6: Symmetry}

\sectionheading{Sections 7--9: Group Operations and the Counting Formula}

\begin{exercise}{S.6 (Supplementary)}
\begin{enumerate}[label=\textbf{(\alph*)}]
  \item Let $p$ be a prime and let $G$ be a group of order $p^k$ acting
        on a finite set $S$. Prove that if $p$ does not divide $|S|$, then
        the action has a \emph{fixed point}, that is, some $s\in S$ with
        $g\cdot s=s$ for every $g\in G$.
  \item Prove that every action of a group of order $15$ on a set with
        $7$ elements has a fixed point.
\end{enumerate}
\end{exercise}

\begin{exercise}{9.2 (Adapted)}
Let $G$ be the group of rotational symmetries of a cube, and let $V$, $E$,
$F$ denote the sets of vertices, edges, and faces of the cube. Fix a
vertex $v$, an edge $e$, and a face $f$, and let $G_v$, $G_e$, $G_f$ be
their stabilizers.
\begin{enumerate}[label=\textbf{(\alph*)}]
  \item Determine $|G_v|$, $|G_e|$, and $|G_f|$, and describe each
        stabilizer geometrically.
  \item For each of the three subgroups $G_v$, $G_e$, $G_f$, decompose each of
        $V$, $E$, $F$ into orbits. Record your answers as sums
        of orbit sizes (for example $8=1+1+3+3$).
\end{enumerate}
\end{exercise}

\sectionheading{Section 12: Finite Subgroups of the Rotation Group}

\begin{exercise}{S.7 (Supplementary)}
\begin{enumerate}[label=\textbf{(\alph*)}]
  \item Use Euler's rotation theorem to prove that the composition of two rotations
        of $\R^3$ (about axes through the origin) is again a rotation.
  \item Let $\rho_1$ be the rotation by $\pi/2$ about the $e_3$-axis and
        $\rho_2$ the rotation by $\pi/2$ about the $e_1$-axis, each
        counterclockwise when viewed from the tip of the axis vector.
        Find a pole and the angle of rotation of $\rho_2\circ\rho_1$.
        \textit{Hint.} The trace of a rotation through angle $\theta$
        is $1+2\cos\theta$. Why?
\end{enumerate}
\end{exercise}

\begin{exercise}{S.8 (Supplementary)}
For a finite group $G\leq\SO_3(\R)$ of order $N$, let $\mathcal P$ be its
set of poles: the unit vectors fixed by some nonidentity element of $G$.
Recall the pole-counting formula
\[
  \sum_{p\in\mathcal P}\bigl(|G_p|-1\bigr)=2(N-1).
\]
\begin{enumerate}[label=\textbf{(\alph*)}]
  \item Let $O$ be the group of rotational symmetries of the cube
        $[-1,1]^3$, which has order $24$. Describe all the poles of $O$, sort them into
        $O$-orbits, give the order of the stabilizer of a pole in each orbit,
        and verify the formula. Which row of the table from
        Lecture 24 does this realize?
  \item Let $n\geq3$, and let $D_n\leq\SO_3(\R)$ be the group of rotational symmetries
        of a regular $n$-gon in the $xy$-plane centered at the origin.
        It consists of $n$ rotations about the $z$-axis and $n$ half-turns
        about axes in the $xy$-plane. Do the same for $D_n$, and check that you obtain the
        dihedral pattern $(r_1,r_2,r_3)=(2,2,n)$.
\end{enumerate}
\end{exercise}

\begin{exercise}{12.3 (Adapted)}
Let $O$ be the group of rotational symmetries of a cube. Its elements may
exchange the two endpoints of a long diagonal. Determine the stabilizer in
$O$ of a long diagonal (a line through two opposite vertices), list its
elements geometrically, and identify it with a familiar group.
\end{exercise}

\vfill
{\small\color{darkgray}Compiled \today.}

\end{document}
